\documentclass[11pt,a4paper]{article}
% =====================================================================
%  Shared preamble for the Part V lecture notes (lectures 19-22).
%
%  These four lectures sit outside Geron, so the notes ARE the primary
%  source and are examined on the same terms as the textbook parts.
%  Everything here is chosen to make that readable on paper: one column,
%  generous measure, and the same six-colour palette as the slides so a
%  student moving between the two is not re-learning what red means.
%
%  Build with pdflatex (twice, for the toc and the refs):
%      cd notes && latexmk -pdf lecture-19.tex
%  or  make -C notes
% =====================================================================
\usepackage[T1]{fontenc}
\usepackage[utf8]{inputenc}
\usepackage{newtxtext,newtxmath}      % Times-like: prints well, reads well
\usepackage[scaled=0.86]{inconsolata} % code, at a size that sits beside Times
\usepackage{microtype}
% newtx defines \Bbbk, \openbox, \proof and \endproof; amssymb and amsthm
% then redefine them and abort. Freeing the four names before those two
% packages load is the standard fix, and it costs nothing here: we use
% amsthm only for \newtheorem.
\let\Bbbk\relax
\usepackage{amsmath,amssymb}
\let\openbox\relax
\let\proof\relax
\let\endproof\relax
\usepackage{amsthm}
\usepackage{booktabs}
\usepackage{enumitem}
\usepackage{xcolor}
\usepackage{tcolorbox}
\usepackage{titlesec}
\usepackage{graphicx}
\usepackage[hidelinks]{hyperref}
\tcbuselibrary{skins,breakable}

% ---- the course palette, from assets/css/custom.css -------------------
\definecolor{aimlprimary}{HTML}{0B3D62}   % deep blue  - structure
\definecolor{aimlaccent} {HTML}{C0392B}   % brick red  - failures
\definecolor{aimlsuccess}{HTML}{14663A}   % green      - repairs
\definecolor{aimlmath}   {HTML}{6C3483}   % purple     - the derivation
\definecolor{aimlmuted}  {HTML}{4B5563}
\definecolor{aimlrule}   {HTML}{B0BCC7}
\definecolor{aimlsoft}   {HTML}{F4F7F9}

\usepackage[a4paper,margin=32mm,top=30mm,bottom=30mm]{geometry}
\setlength{\parskip}{0.45\baselineskip}
\setlength{\parindent}{0pt}
\linespread{1.06}

% ---- headings --------------------------------------------------------
\titleformat{\section}{\Large\bfseries\color{aimlprimary}}{\thesection}{0.7em}{}
\titleformat{\subsection}{\large\bfseries\color{aimlprimary}}{\thesubsection}{0.7em}{}
\titleformat{\subsubsection}{\bfseries\color{aimlmuted}}{}{0em}{}
\titlespacing*{\section}{0pt}{1.6\baselineskip}{0.5\baselineskip}
\titlespacing*{\subsection}{0pt}{1.2\baselineskip}{0.35\baselineskip}

% ---- boxes -----------------------------------------------------------
% One box per role, and the role is the colour. A student who has seen the
% slides already knows what each one means before reading a word of it.
\newtcolorbox{keybox}[1][]{
  breakable, enhanced, colback=aimlsoft, colframe=aimlprimary,
  boxrule=0.4pt, left=8pt, right=8pt, top=6pt, bottom=6pt,
  fonttitle=\bfseries\small, coltitle=white, colbacktitle=aimlprimary,
  attach boxed title to top left={yshift=-2mm, xshift=4mm},
  boxed title style={boxrule=0pt, arc=1pt}, #1}

\newtcolorbox{failbox}[1][]{
  breakable, enhanced, colback=white, colframe=aimlaccent,
  boxrule=0.9pt, left=8pt, right=8pt, top=6pt, bottom=6pt,
  fonttitle=\bfseries\small, coltitle=white, colbacktitle=aimlaccent,
  attach boxed title to top left={yshift=-2mm, xshift=4mm},
  boxed title style={boxrule=0pt, arc=1pt}, #1}

\newtcolorbox{derivbox}[1][]{
  breakable, enhanced, colback=white, colframe=aimlmath,
  boxrule=0.6pt, left=8pt, right=8pt, top=6pt, bottom=6pt,
  fonttitle=\bfseries\small, coltitle=white, colbacktitle=aimlmath,
  attach boxed title to top left={yshift=-2mm, xshift=4mm},
  boxed title style={boxrule=0pt, arc=1pt}, #1}

% An exercise. Deliberately the same shape as the notebook's `try` field:
% one change, and what should happen to the output.
\newtcolorbox{trybox}[1][]{
  breakable, enhanced, colback=white, colframe=aimlsuccess,
  boxrule=0.5pt, left=8pt, right=8pt, top=6pt, bottom=6pt,
  fonttitle=\bfseries\small, coltitle=white, colbacktitle=aimlsuccess,
  attach boxed title to top left={yshift=-2mm, xshift=4mm},
  boxed title style={boxrule=0pt, arc=1pt}, #1}

% ---- small semantics -------------------------------------------------
\newcommand{\scope}[1]{\textcolor{aimlmuted}{\small #1}}
\newcommand{\term}[1]{\textbf{#1}}
\newcommand{\code}[1]{\texttt{\small #1}}
\newcommand{\lecture}[1]{Lecture~#1}

\newtheorem{definition}{Definition}[section]
\newtheorem{proposition}[definition]{Proposition}

% ---- title block -----------------------------------------------------
% \notestitle{19}{Information retrieval: the lexical foundation}{SciFact (BEIR)}
\newcommand{\notestitle}[3]{%
  \thispagestyle{empty}
  \begin{flushleft}
    {\color{aimlmuted}\small\sffamily
     APPLICATIONS OF MACHINE LEARNING \quad$\cdot$\quad
     BSc Mathematics of Artificial Intelligence}\\[2mm]
    {\color{aimlrule}\rule{\textwidth}{0.8pt}}\\[4mm]
    {\color{aimlmuted}\large Lecture #1 \quad$\cdot$\quad Part V \quad$\cdot$\quad
     Extended notes}\\[3mm]
    {\Huge\bfseries\color{aimlprimary}\baselineskip=1.15\baselineskip #2\par}\vspace{4mm}
    {\color{aimlmuted}\large Dataset: #3}\\[3mm]
    {\color{aimlrule}\rule{\textwidth}{0.8pt}}
  \end{flushleft}
  \vspace{2mm}
  \begin{keybox}[title={Why these notes exist}]
    Lectures 19--22 sit outside G\'eron, the course's primary text. For those
    four lectures \emph{these notes are the primary source}, and they are
    examined on exactly the same terms as the textbook parts. Everything in
    the slide deck for this lecture is here, with the arguments written out at
    length rather than compressed onto a slide; the numbers are the ones the
    accompanying notebook prints, and every one of them is a measurement on
    the corpus named above rather than a figure quoted from a paper.
  \end{keybox}
  \vspace{1mm}
  {\small\tableofcontents}
  \vspace{2mm}
  \noindent{\color{aimlrule}\rule{\textwidth}{0.4pt}}
  \vspace{1mm}
}


\begin{document}
\notestitlefor{12}{III}{Convolutional networks}{Flowers102}{Chapter 12}

\section{What breaks when the image gets bigger}

Three lectures of neural networks, on flattened greyscale images. Flattening a
224 by 224 colour image gives 150{,}528 inputs, and a single dense layer of
1{,}000 units on that is 150 million parameters --- for one layer, on one image
size. Every one has to be learned from data, and none of them knows that two
adjacent pixels are related.

Today's method fixes both problems at once, and the derivation says exactly how
much it saves.

\subsection{The brief}

A horticultural wholesaler receives mixed pallets of cut stems and must sort
them by species before the auction floor. One overhead camera per lane, one
photograph per stem, and a human inspector who is the bottleneck.

In their words: \emph{``Tell us which of our species it is. If you cannot tell,
say so and we will look at it ourselves --- but do not tell us the wrong
one.''} Two requirements hiding in one sentence: a classifier over a fixed
list, and a usable notion of confidence. We build the first today.

\begin{center}
\small
\begin{tabular}{lll}
\toprule
& \textbf{Fashion-MNIST} & \textbf{Flowers102} \\
\midrule
Image & 28 by 28, grey & variable, colour \\
Classes & 10 & 102 \\
Training images & 60{,}000 & 1{,}020 \\
Images per class & 6{,}000 & 10 \\
Subject fills the frame & always & sometimes \\
\bottomrule
\end{tabular}
\end{center}

Ten times the classes and 59 times fewer labelled images. The split shipped
with the dataset is 1{,}020 train, 1{,}020 val and 6{,}149 test --- six times as
many test images as training images, and the training split is balanced at
exactly ten per class where the test split is not.

\begin{keybox}[title={The standing constraint of applied vision}]
Images are cheap and labels are not. Somebody had to look at every photograph
and name the species. Every application in this part of the course has more
pixels than it can use and fewer labels than it wants --- so a method that
needs thousands of labelled examples per class is not a method you can deploy
here.
\end{keybox}

What is hard here is position (the flower is not centred and not at a fixed
scale), background, between-class similarity (several species separated by
petal count and little else), and within-class variation. \emph{Only the first
has a structural answer}, and it is the one this chapter is about.

\section{Derivation: weight sharing, equivariance, and memory}

You replace a matrix multiplication with a convolution and the network becomes
trainable. What exactly did you buy, and what did you give up? Three answers: a
parameter count, a symmetry, and a memory bill that is not where anyone expects
it.

\subsection{The parameter count, on the same image}

Fix the input and output shapes so the comparison is between two ways of
computing the same sized thing: input 3 by 128 by 128, which is 49{,}152
values, and output 32 by 128 by 128, which is 524{,}288.

A dense layer connects every output to every input,
\[ \#\text{weights} = n_{\text{in}}n_{\text{out}} = (3HW)(32HW) = 96H^2W^2, \]
which is 25{,}769{,}803{,}776 weights --- 96 GB of float32 for one layer,
before any gradient or optimiser state.

A convolution connects every output to a $k$ by $k$ window of every input
channel and \emph{uses the same weights at every position},
\[ \#\text{weights} = C_{\text{out}}C_{\text{in}}k^2 = 32\cdot3\cdot7\cdot7, \]
which is 4{,}704. Note that $H$ and $W$ do not appear: the image size is
genuinely absent from the count. The ratio is $H^2W^2/k^2$, which here is
5{,}478{,}275 times fewer weights for the same output shape --- and with the 3
by 3 filters used everywhere else in the network, 29{,}826{,}162.

\begin{failbox}[title={What the saving is not}]
It is \emph{not} a saving in arithmetic. The convolution still computes
$C_{\text{out}}HW$ outputs, each a sum over $C_{\text{in}}k^2$ terms --- the
same order of multiplications as the dense layer's, to within the window size.

What shrank is the number of \emph{distinct numbers that have to be stored and
learned}. The win is in memory and in statistics, not in floating-point
operations, and confusing the two is how people end up expecting convolutions
to be fast.
\end{failbox}

\begin{keybox}[title={Weight sharing is a constraint, not a trick}]
A dense layer can represent any convolution: set its weight matrix to the
block-Toeplitz matrix that repeats the kernel. Nothing is lost in expressive
power.

The convolution is that matrix with its parameters \emph{tied}. Tying them says
a pattern means the same thing wherever it appears --- true of photographs,
false of a spreadsheet. This is Lecture 5's bias--variance trade made
architecturally: a large reduction in variance, bought with an assumption we
are prepared to defend.
\end{keybox}

\subsection{What that assumption is, precisely}

Let $T_{\mathbf{v}}$ shift an image by $\mathbf{v}$, so
$(T_{\mathbf{v}}x)(\mathbf{p}) = x(\mathbf{p} - \mathbf{v})$. A map $f$ is
\term{equivariant} to translation when
$f(T_{\mathbf{v}}x) = T_{\mathbf{v}}f(x)$ for every $\mathbf{v}$: shift then
compute, or compute then shift, is the same answer.

\begin{derivbox}[title={Convolution is equivariant --- one change of variable}]
Write $(w * x)(\mathbf{p}) = \sum_{\mathbf{u}} w(\mathbf{u})
x(\mathbf{p}+\mathbf{u})$. Then
\[
\begin{aligned}
(w * T_{\mathbf{v}}x)(\mathbf{p})
 &= \sum_{\mathbf{u}} w(\mathbf{u})\,x(\mathbf{p}+\mathbf{u}-\mathbf{v}) \\
 &= (w * x)(\mathbf{p}-\mathbf{v})
  = \bigl(T_{\mathbf{v}}(w * x)\bigr)(\mathbf{p}).
\end{aligned}
\]
The proof works because $w$ does not depend on $\mathbf{p}$ --- which is
exactly what weight sharing is. Untie the weights and the middle equality
fails.
\end{derivbox}

The equality is not exact everywhere. At the border, zero padding invents input
that was not there, so a shift moves real content into invented content. Under
stride greater than one, the output grid samples every $s$-th position, so
equivariance survives only for shifts that are multiples of $s$. And in
floating point, if the two computations sum the same terms in different orders.
All three are why the next paragraph is a measurement rather than an assertion.

Shifting a trained layer's input by 16 pixels and dropping a 24-pixel border,
the largest absolute difference on the interior is \textbf{0}, against a
largest activation of 4.15 in the same region. Bit-for-bit identical, not
merely close --- on this backend the two computations perform the same
additions in the same order. Check it rather than assuming it, because on
another backend it will not be exactly zero.

\subsection{Equivariance is not invariance}

A map is \term{invariant} when $g(T_{\mathbf{v}}x) = g(x)$: under a shift the
output does not change, and the position is not kept. Invariance is strictly
the stronger and strictly the lossier of the two, and it is what you get by
discarding the equivariant structure.

Pooling manufactures it. A 2 by 2 max pool is invariant to any shift that keeps
the maximum inside its own window; stack four stages and the tolerated shift
grows to about eight pixels in the last feature map; a global pool over the
whole map is invariant to every shift, because there is no position left to be
sensitive to. The invariance is approximate and it is \emph{built}, one factor
of two at a time --- nothing in the architecture declares it.

\begin{center}
\small
\begin{tabular}{lr}
\toprule
\textbf{Cosine similarity, image against image shifted 16 pixels} & \\
\midrule
the spatial feature maps, flattened & 0.6379 \\
after a global max pool & 1.0000 \\
\bottomrule
\end{tabular}
\end{center}

The map changed by 36.2\%; the pooled vector changed by 0.0004\%.

Classification wants the second line: ``this is a sunflower'' is true wherever
the sunflower is, so a representation that still carries the position is
carrying a nuisance variable, and discarding it reduces variance at no cost in
bias \emph{provided the assumption really holds}.

\begin{failbox}[title={Failure condition --- why per-pixel prediction cannot afford it}]
Segmentation asks, for every pixel, which class is it --- and the answer
\emph{is} the position. Our input grid is 128 by 128 and the last feature map
is 8 by 8, so one cell answers for 256 pixels: a 256-fold loss of spatial
resolution, and two flower boundaries 15 pixels apart are the same cell.

Lecture 14 has to put that resolution back, and the whole of its architecture
is about how.
\end{failbox}

\subsection{Where the memory goes}

Your network has 4{,}807{,}494 parameters and your session died with
\emph{out of memory}. Which of those two facts caused the other?

\begin{center}
\small
\begin{tabular}{lrr}
\toprule
\textbf{What} & \textbf{Bytes each} & \textbf{Total} \\
\midrule
weights & 4 & 18.3 MB \\
gradients, one per weight & 4 & 18.3 MB \\
Adam's first moment & 4 & 18.3 MB \\
Adam's second moment & 4 & 18.3 MB \\
\midrule
everything parameter-shaped & 16 & 73.4 MB \\
\bottomrule
\end{tabular}
\end{center}

Sixteen bytes per parameter, not four --- the factor people forget, and still
not the answer.

Lecture 10's derivation, in bytes: the reverse sweep evaluates each Jacobian
\emph{at the value the forward pass reached}, so every layer's output stays in
memory from the moment it is computed until the backward pass consumes it. That
is not an implementation detail of PyTorch; it is what reverse-mode automatic
differentiation \emph{is}.

Counted for the actual network, the stored activations are 5{,}964{,}646 floats
per image --- 22.8 MB. At batch 32 that is 728 MB: 9.9 times everything
parameter-shaped put together, and about forty times the weights alone.

And it is not spread evenly. The activation cost of a layer is $C\times H\times
W$, and channels double as the map quarters, so the total \emph{halves} at
every pooling stage: the first convolution's output is 64 MB at batch 32, the
third 32 MB, the last 8 MB. 49\% of the convolutional activation memory is in
the first two layers --- the layers with almost no parameters.

\begin{keybox}[title={Parameters and activations live at opposite ends}]
Activations concentrate in the \emph{early} layers, because $H$ and $W$ are
still large. Parameters concentrate in the \emph{late} layers, because
$C_{\text{in}}C_{\text{out}}$ is large and $k$ is fixed.

Two different bills, sent by two different halves of the same network --- which
is why ``make the model smaller'' and ``make the run fit in memory'' are
different problems with different answers.
\end{keybox}

Checked against the allocator, the prediction by counting outputs is 728 MB and
the measurement is 568. The measurement is \emph{lower} because the allocator
frees an activation as soon as nothing can still need it, and in a plain
\code{Sequential} the ReLU outputs of already-consumed layers qualify. The
prediction is the ceiling, which is the number you want when sizing a run.

\begin{keybox}[title={When a run runs out of memory}]
Halve the batch, not the model. Activation memory is linear in the batch size
and the parameter memory does not move at all, so the first lever costs you
nothing but wall clock. In order: smaller batch, then smaller input resolution
(quadratic), then gradient checkpointing, then mixed precision.

Reducing the parameter count is near the bottom of that list --- and on this
network it would mostly delete the dense head, which is not where the memory
is.
\end{keybox}

\section{Convolution, pooling, and an architecture}

A feature worth detecting in the top left corner is worth detecting in the
bottom right, so use the same weights everywhere. That is the whole idea.

For an input with $C$ channels and a kernel of size $k$ by $k$,
\[ z^{(f)}_{ij} = b^{(f)} + \sum_{c=1}^{C}\sum_{u=0}^{k-1}\sum_{v=0}^{k-1}
   w^{(f)}_{c,u,v}\,x_{c,\,i+u,\,j+v} \]
--- a dot product between the kernel and a window, repeated at every position,
with the sum over $c$ being why a filter reads all input channels at once.

Three numbers describe it: the kernel size $k$, how much of the input one
output sees; the stride $s$, how far the window moves; and the padding $p$, how
many zeros ring the input. Then
\[ H_{\text{out}} = \left\lfloor\frac{H_{\text{in}} + 2p - k}{s}\right\rfloor
   + 1, \]
and with $s = 1$ and $p = \lfloor k/2\rfloor$ the output keeps the input's
height and width.

\begin{failbox}[title={Failure condition --- the default nobody asked for}]
\code{padding} defaults to \emph{zero}, which silently shrinks every feature
map by $k-1$ and will eventually give you a shape error twelve layers later.
Reviewer question 5, every time.
\end{failbox}

Pooling halves the height and width by keeping the largest value in each 2 by 2
block. Three quarters of the activations disappear, and with them three
quarters of the memory and arithmetic in every later layer; it has no
parameters; and each later unit sees a larger piece of the original image. The
fourth consequence is the invariance the derivation spent twenty minutes on.

Every convolution is followed by a batch normalisation, because ten layers of
untreated convolutions on 1{,}020 images do not train --- the loss sits at
$\log 102 \approx 4.62$ and the network predicts one class. That is Lecture
11's $\rho \ne 1$ in a new architecture, and here, unlike there, initialisation
alone does not reach it. Lecture 11 measured batch normalisation making a
repaired network \emph{worse} by 4.3 points; the size of an effect depends on
what else is in the system.

\begin{center}
\small
\begin{tabular}{llr}
\toprule
\textbf{Stage} & \textbf{Layers} & \textbf{Output} \\
\midrule
1 & conv 7 by 7 $\times$32, conv 3 by 3 $\times$32, pool & 32 by 64 by 64 \\
2 & conv $\times$64, conv $\times$64, pool & 64 by 32 by 32 \\
3 & conv $\times$128, conv $\times$128, pool & 128 by 16 by 16 \\
4 & conv $\times$256, pool & 256 by 8 by 8 \\
head & flatten, dense 256, dropout, dense 102 & 102 \\
\bottomrule
\end{tabular}
\end{center}

\begin{center}
\small
\begin{tabular}{lrr}
\toprule
\textbf{Part} & \textbf{Parameters} & \textbf{Share} \\
\midrule
Every convolution and batch-norm layer together & 586{,}720 & 12.2\% \\
One dense layer, 16{,}384 to 256 & 4{,}194{,}560 & 87.3\% \\
Output layer, 256 to 102 & 26{,}214 & 0.5\% \\
\midrule
Total & 4{,}807{,}494 & 100\% \\
\bottomrule
\end{tabular}
\end{center}

Nine parameters in ten are in the layer that is \emph{not} convolutional. And
4{,}807{,}494 parameters against 1{,}020 training images is 4713 parameters per
training image: there is nothing in the architecture stopping this network from
storing the training set, and dropout at 0.5 is the only thing asked to prevent
it.

The training loop is exactly Lecture 10's five lines. Convolution changed the
model; it did not change the loop, because autograd does not care what the
modules are.

\section{What it costs, and what it cannot do}

\begin{center}
\small
\begin{tabular}{lrr}
\toprule
\textbf{Learning rate} & \textbf{Mean test accuracy} & \textbf{sd over seeds} \\
\midrule
0.001, Adam's default & 2.7\% & 1.31 pts \\
0.0003 & 14.7\% & 2.59 pts \\
\bottomrule
\end{tabular}
\end{center}

12.1 points, from one keyword argument nobody wrote down. The default does not
diverge --- it plateaus low, which is the failure mode that does not announce
itself.

Eighty epochs at batch 32 is 32 steps per epoch and 2{,}560 optimiser steps in
total, 253 seconds for the seed-42 run. The training curve reaches 71.8\% while
the validation curve stops climbing early: 58 points of the training accuracy
is memorisation, which is the shape Lecture 5 named as variance. The last two
thirds of the wall clock bought almost nothing.

\begin{center}
\small
\begin{tabular}{lrr}
\toprule
\textbf{Model} & \textbf{Test accuracy} & \textbf{Wall clock} \\
\midrule
Uniform random guess & 0.98\% & --- \\
Always the commonest species & 3.87\% & --- \\
Convolutional network, from scratch & 15.3\% & 364 s \\
\midrule
What the operator needs & 90\% & --- \\
\bottomrule
\end{tabular}
\end{center}

Four times the majority baseline, and a long way from deployable.

\begin{keybox}[title={Why five runs and not one}]
Seeds 42 to 46 give 11.0\%, 15.7\%, 18.1\%, 16.1\% and 15.7\% --- 7.1 points
between best and worst with nothing changed but the seed.

One run of a network this size on 1{,}020 images is an anecdote, and any
comparison you make from here that is smaller than 7.1 points is not a
comparison.
\end{keybox}

\subsection{Look at what it learned}

The first layer's weights are 4{,}704 numbers as 32 filters of 7 by 7 by 3 ---
small enough to look at. Against the initialisation they are still almost the
same vector --- the mean absolute change is 23.9\% of the mean weight, and 247
of the 4{,}704 changed sign.

The weights moved. \emph{No structure appeared} --- and 253 seconds of training
bought that. Two reasons, both worth knowing. The gradient reaching this layer
is small, because ten layers of chain rule separate it from the loss and 1{,}020
images give only 2{,}560 optimiser steps. And its \emph{scale does not matter}:
a batch norm follows immediately, so normalising the output makes the overall
size of these weights irrelevant to the function, and there is no gradient
pressure on it at all.

The textbook picture of a first layer --- colour blobs, oriented edges --- is
real, and you will see it next lecture. It comes from a corpus a thousand times
this size.

\section{The notebook}

\begin{trybox}[title={What to change}]
Halve the batch from 32 to 16 and watch the allocator. The parameter figure
does not move and the activation figure halves, which is the memory argument
measured rather than read --- and it is also the first lever the rule above
tells you to pull.
\end{trybox}

\section{Exercises}

\begin{enumerate}[leftmargin=1.6em,itemsep=4pt]
  \item A dense layer and a convolution produce the same output shape from the
        same image. Give both parameter counts symbolically, and state which
        quantity is absent from the second. \hfill \scope{[5]}
  \item Prove that convolution is equivariant to translation, and name the
        property of the kernel the proof depends on. \hfill \scope{[4]}
  \item Explain why pooling's invariance helps classification and is fatal for
        per-pixel prediction. \hfill \scope{[4]}
  \item A training run fails with out-of-memory. The model has 4.8 million
        parameters. State which of parameters or activations you would
        investigate first, with the reason, and give the first lever you would
        pull. \hfill \scope{[5]}
  \item A first layer's filters look like noise after training. Give two
        reasons this can happen even when the network as a whole is learning.
        \hfill \scope{[4]}
\end{enumerate}

\section*{Summary}
\addcontentsline{toc}{section}{Summary}

Weight sharing costs 5{,}478{,}275-fold fewer parameters than the dense layer
with the same output shape --- a saving in memory and in statistics, not in
arithmetic. It is precisely the condition for translation equivariance, and the
proof is one change of variable that depends on the kernel not varying with
position. Pooling trades that equivariance for invariance, which classification
wants and per-pixel prediction cannot have: measured, the spatial maps of a
shifted image agree at cosine 0.6379 and the globally pooled vectors at 1.0000.

The memory bill is not where anyone expects. Parameters and their optimiser
state are 73.4 MB at sixteen bytes each; activations at batch 32 are 728 MB,
concentrated in the early layers where there are almost no parameters. So halve
the batch, not the model.

Trained from scratch on 1{,}020 images the network reaches 15.3\% against a
3.87\% majority baseline, with 7.1 points between the best and worst of five
seeds. Its learning rate is worth 12.1 points, its training accuracy reaches
71.8\% with 58 points of that being memorisation, and its first-layer filters
are still noise --- for two reasons, one of which is that a batch norm
downstream removes any gradient pressure on their scale.

\end{document}
